\begingroup
\begin{table}[H]
\centering
\caption{Regional peak-coincidence factors used to convert average AI load to peak-load-equivalent demand.}
\label{tab:peak-coincidence-gamma}
\begin{tabular}{lllll}
\toprule
Grid region & Scenario & Period & $\Gamma$ & Three-hour peak window \\
\midrule
PJM & low/mid/high & 2025-2035 & 1.015 & yes \\
MISO & low/mid/high & 2025-2035 & 1.019 & yes \\
SPP & low/mid/high & 2025-2035 & 1.019 & yes \\
ERCOT & low/mid/high & 2025-2035 & 1.019 & yes \\
CAISO & low/mid/high & 2025-2035 & 1.004 & yes \\
NYISO & low/mid/high & 2025-2035 & 1.009 & yes \\
ISO-NE & low/mid/high & 2025-2035 & 1.009 & yes \\
\bottomrule
\end{tabular}
\tabnote{The peak-coincidence factor converts allocated average AI demand into the peak-load-equivalent quantity used in the annual margin loop. Counties outside the seven analyzed grid regions remain in the national spatial inventory but do not receive a common ISO peak-window adjustment.}
\end{table}
\endgroup
